%%==========================================================================
%%  Research Statement — Zheng Ma
%%==========================================================================
\documentclass[11pt, a4paper]{article}

%--- Encoding & fonts -------------------------------------------------------
\usepackage[T1]{fontenc}
\usepackage[utf8]{inputenc}
\usepackage{palatino}
\usepackage{microtype}

%--- Page geometry ----------------------------------------------------------
\usepackage[
  top=2.4cm, bottom=2.4cm,
  left=2.8cm, right=2.8cm,
  headheight=14pt
]{geometry}

%--- Color (minimal — dark ink only, subdued link color) --------------------
\usepackage{xcolor}
\definecolor{linkblue}{HTML}{1A3A5C}

%--- Math -------------------------------------------------------------------
\usepackage{amsmath, amssymb, amsthm}

%--- Headers / footers ------------------------------------------------------
\usepackage{fancyhdr}
\pagestyle{fancy}
\fancyhf{}
\renewcommand{\headrulewidth}{0.4pt}
\fancyhead[L]{\small\textit{Research Statement \,---\, Zheng Ma}}
\fancyhead[R]{\small\textit{April 2026}}
\fancyfoot[C]{\small\thepage}

%--- Section titles (traditional academic) ----------------------------------
\usepackage{titlesec}
\titleformat{\section}
{\large\bfseries}{}{0em}{}[\vspace{-0.5ex}\rule{\linewidth}{0.5pt}\vspace{0.3ex}]
\titlespacing*{\section}{0pt}{2.6ex plus 0.5ex minus 0.3ex}{0.6ex plus 0.1ex}
\titleformat{\subsection}
{\normalsize\bfseries\itshape}{}{0em}{}
\titlespacing*{\subsection}{0pt}{1.4ex plus 0.3ex minus 0.2ex}{0.2ex}

%--- Lists ------------------------------------------------------------------
\usepackage{enumitem}
\setlist[itemize]{leftmargin=1.6em, itemsep=0.1em, topsep=0.4em, parsep=0em}
\setlist[enumerate]{leftmargin=1.6em, itemsep=0.1em, topsep=0.4em, parsep=0em}

%--- Hyperref ---------------------------------------------------------------
\usepackage[
  colorlinks=true,
  linkcolor=linkblue,
  citecolor=linkblue,
  urlcolor=linkblue,
  pdftitle={Research Statement --- Zheng Ma},
  pdfauthor={Zheng Ma}
]{hyperref}

%--- Bibliography -----------------------------------------------------------
\usepackage[
  backend=biber,
  style=numeric-comp,
  sorting=none,
  maxnames=6,
  giveninits=true,
  doi=false,
  isbn=false,
  url=false
]{biblatex}
\addbibresource{references.bib}

%--- Typography -------------------------------------------------------------
\setlength{\parindent}{1.5em}
\setlength{\parskip}{0pt}
\frenchspacing
\usepackage{xurl}   % allow line breaks in URLs

%%==========================================================================
%%  DOCUMENT
%%==========================================================================
\begin{document}

%%--------------------------------------------------------------------------
%%  TITLE BLOCK
%%--------------------------------------------------------------------------
\begin{center}
  {\LARGE\bfseries Research Statement}\\[0.8em]
  {\large Zheng Ma}\\[0.3em]
  {\normalsize
    Tenure-Track Associate Professor of Mathematics,
  Shanghai Jiao Tong University}\\[0.25em]
  {\normalsize
    \href{mailto:zhengma@sjtu.edu.cn}{\texttt{zhengma@sjtu.edu.cn}}
    \enspace\textbar\enspace
    \href{https://zheng-home.netlify.app}{\texttt{zheng-home.netlify.app}}
    \enspace\textbar\enspace
    \href{https://github.com/mazhengcn}{\texttt{github.com/mazhengcn}}
  }
\end{center}
\medskip
\noindent\rule{\linewidth}{1pt}

\bigskip

%%--------------------------------------------------------------------------
%%  SECTION 1: INTRODUCTION AND VISION
%%--------------------------------------------------------------------------
\section{Introduction and Research Vision}

We stand at a remarkable inflection point in computational science.  The most
pressing open problems of our time---controlled nuclear fusion, precision
subsurface imaging for climate-critical carbon sequestration, next-generation
semiconductor design, and predictive modeling of radiative transport in
astrophysical and biomedical settings---share a common mathematical core: they
require solving kinetic equations and associated inverse problems at scales,
speeds, and levels of physical fidelity that are fundamentally inaccessible to
classical numerical methods.

Kinetic theory provides the most complete description of matter at the
mesoscale, the regime bridging quantum mechanics and macroscopic fluid
dynamics.  The Boltzmann equation, the radiative transfer equation (RTE), and
the Vlasov--Maxwell system underpin models of plasma confinement in tokamaks,
photon transport in stellar atmospheres and radiation oncology planning, and
carrier dynamics in semiconductor devices.  These equations define
distributions over phase space of dimension $2d$ ($d = 3$ in three spatial
dimensions), rendering direct discretization intractable for realistic
problems.  Compounding the curse of dimensionality is the ubiquitous presence
of multiple temporal and spatial scales---from the mean free path of a single
particle collision to macroscopic fluid motion---that force classical solvers
to use extremely fine computational grids even when the macroscopic solution
is globally smooth.

Deep learning has demonstrated a striking capacity to approximate
high-dimensional functions and solution operators with sample complexity far
better than classical worst-case theory would predict.  Yet the na\"{i}ve
application of neural networks to physics problems is beset by fundamental
limitations: models trained without respect for physical structure violate
conservation laws, behave catastrophically in stiff multiscale regimes where
the Knudsen number $\varepsilon \ll 1$, and offer no theoretical guarantees
for out-of-distribution generalization.

My research program addresses these limitations through what I call
\emph{full-stack AI for Science}: a vertically integrated methodology that
combines (i)~rigorous mathematical analysis of the underlying physics and the
learning algorithm, (ii)~neural network architectures that encode physical
structure as \emph{hard constraints} rather than soft regularization
penalties, (iii)~scalable implementation exploiting modern accelerated
hardware, and (iv)~production-quality open-source software that makes these
methods accessible to the broader scientific community.  This program has
produced more than 30 peer-reviewed articles and a suite of open-source
packages---hosted at
\href{https://github.com/mazhengcn}{\texttt{github.com/mazhengcn}}---that
are actively used by research groups in applied mathematics, physics, and
engineering.  The work spans four interconnected thrusts: \textbf{(1)
structure-preserving AI for multiscale kinetic equations}; \textbf{(2) a
mathematical theory of deep neural network training}; \textbf{(3)
AI-driven methods for PDE-constrained inverse problems}; and \textbf{(4)
fast spectral algorithms} that serve as both the theoretical bedrock and the
benchmark gold standard for the AI-based methods.

%%--------------------------------------------------------------------------
%%  SECTION 2: KINETIC EQUATIONS
%%--------------------------------------------------------------------------
\section{Structure-Preserving AI for Multiscale Kinetic Equations}

\noindent\textit{Why it matters.}\enspace
Kinetic equations are among the most computationally expensive models in all
of science.  A single RTE simulation in nuclear reactor transport involves
$10^8$--$10^{12}$ phase-space unknowns, with run times on classical
supercomputers measured in hours.  Fusion plasma simulations face a starker
challenge: real-time feedback control of a tokamak plasma requires millions
of Vlasov--Maxwell evaluations per second, a throughput that is five to six
orders of magnitude beyond what current particle-in-cell (PIC) codes can
deliver.  Equally important is the \emph{multiscale} structure of these
equations: in the limit $\varepsilon \to 0$ (small Knudsen number), the
kinetic model degenerates to a macroscopic diffusion or fluid system, but
classical numerical discretizations become stiff and impose a time-step
restriction $\Delta t = O(\varepsilon)$, adding an entire order of
computational difficulty to an already intractable problem.

\subsection{Asymptotic-Preserving Neural Networks}

The central contribution of my group in this area is the development of
\emph{Asymptotic-Preserving Neural Networks}
(APNNs)~\cite{APNN2023,APNN2024}: a fundamentally new class of
physics-informed neural network architectures that provably inherit the
asymptotic-preserving (AP) property of the best classical schemes.  The
essential novelty is architectural.  We encode the micro--macro decomposition
$f = \rho M + \varepsilon g$ (where $M$ is the local Maxwellian equilibrium
and $g$ is the microscopic deviation) directly into the network's structural
inductive bias, rather than including it as a soft loss-function term that
can be violated during optimization.  The ``macro'' branch of the network is
constrained to satisfy the correct limiting diffusion equation as
$\varepsilon \to 0$, and the ``micro'' branch is constrained to vanish in the
same limit.  This is, to our knowledge, the first class of neural network
methods for kinetic equations with a rigorous, verifiable AP guarantee
independent of training quality or random seed.

The framework has been applied across the canonical multiscale kinetic
systems:
\begin{itemize}
  \item \emph{Linear transport equations}~\cite{APNN2023,APNN2024}: uniform
    accuracy from the kinetic regime to the diffusion limit without mesh
    refinement, on problems where standard PINNs fail completely for
    $\varepsilon < 10^{-3}$.

  \item \emph{Vlasov--Poisson--Fokker--Planck system in the high-field
    regime}~\cite{APNN-VPFP2024}: AP behavior in the presence of long-range
    electromagnetic self-interactions, extending the framework to
    nonlinearly coupled systems.

  \item \emph{Gray radiative transfer via even--odd decomposition}
    (AP-CON)~\cite{APCON2024}: operator-learning variants using convolutional
    DeepONets that generalize across $10^4$ distinct optical parameter
    configurations without retraining, achieving accuracy competitive with
    reference $S_N$ solvers while running $10^3\times$ faster at inference.
\end{itemize}

\noindent Reference implementations are publicly available at
\href{https://github.com/mazhengcn/apnns}{\texttt{github.com/mazhengcn/apnns}}
and
\href{https://github.com/mazhengcn/apcons}{\texttt{github.com/mazhengcn/apcons}},
providing training scripts, pre-computed datasets, and documented APIs for
reproducing all published results and adapting the methods to new equations.

\subsection{DeepRTE: A Foundation Model for Radiative Transfer}

A more ambitious development is \textbf{DeepRTE}~\cite{DeepRTE2026}, the
first pre-trained, attention-based \emph{neural operator foundation model}
for the steady-state RTE.  Inspired by the large language model paradigm, the
idea is to train a single model once on a large, systematically generated
dataset of RTE solutions---covering the full range of physically relevant
scattering anisotropy parameters $g \in [0.1, 0.8]$---and then deploy it for
instant inference on any new optical geometry without retraining.  DeepRTE
learns the solution operator $\mathcal{A}: (I^{-};\, \mu_t, \mu_s, p)
\mapsto I$, where $I^{-}$ is the boundary inflow, $\mu_t, \mu_s$ are
attenuation and scattering coefficients, and $p$ is the scattering phase
function, mapping this entire configuration to the full radiation field.

The engineering realization is as important as the scientific contribution.
DeepRTE is implemented in JAX/Flax for hardware-accelerated training and
inference, containerized with Docker for reproducible deployment, and
distributed as a complete package at
\href{https://github.com/mazhengcn/deeprte}{\texttt{github.com/mazhengcn/deeprte}},
with pre-trained model checkpoints and training datasets hosted on Hugging
Face at
\href{https://huggingface.co/mazhengcn/deeprte}{\texttt{mazhengcn/deeprte}}.
The package provides a full training and inference pipeline, enabling
researchers at other institutions to fine-tune the model on domain-specific
geometries---nuclear shielding calculations, photon dosimetry, atmospheric
radiative transfer---without access to our cluster.  This represents a
qualitative shift: rather than a solver for individual instances, DeepRTE is
a \emph{reusable computational asset} that amortizes massive training cost
across the entire scientific community.

\subsection{Asymptotic-Preserving Random Feature Methods}

Beyond neural networks, in collaboration with Jingrun~Chen (USTC), we
developed an \emph{asymptotic-preserving random feature method}
(AP-RFM)~\cite{AP-RFM2025} for multiscale radiative transfer.  Random
feature methods expand the solution in a random basis
$\sigma(\omega_j \cdot x + b_j)$ with fixed random weights, offering
mesh-free approximation with theoretically controllable error at
$O(N)$ complexity.  By encoding the micro--macro decomposition into both
the basis construction and the least-squares fitting, AP-RFM achieves uniform
accuracy across all Knudsen numbers at a cost substantially lower than neural
network training.  This result demonstrates that the AP design principle is
a general mathematical framework---not a property specific to any particular
machine learning architecture.

%%--------------------------------------------------------------------------
%%  SECTION 3: DEEP LEARNING THEORY
%%--------------------------------------------------------------------------
\section{Mathematical Theory of Deep Neural Network Training}

\noindent\textit{Why it matters.}\enspace
Deploying neural networks in safety-critical scientific computing demands
more than empirical validation: it requires a rigorous mathematical
understanding of training dynamics and generalization.  Without such theory,
one cannot predict when a neural network will fail---an unacceptable
situation in nuclear safety transport, clinical imaging, or autonomous
systems.  Beyond safety, a quantitative theory of learning dynamics is an
essential design tool: one must understand how a network's inductive bias
interacts with the spectral properties of the target PDE solution to
guarantee convergence and inform architecture choices.

\subsection{The Frequency Principle}

A central theme of my theoretical work is the \emph{Frequency Principle}
(F-principle)~\cite{FP1}: the finding, first systematically documented by
our group, that gradient descent on neural network loss surfaces fits
low-frequency components of the target function first and progressively
captures higher frequencies.  This is \emph{a priori} surprising from an
approximation-theoretic standpoint, and its mathematical explanation reveals
a deep connection between the spectral properties of the neural tangent
kernel and the implicit regularization induced by early stopping.

In~\cite{FP1} we provided the first systematic Fourier analysis
establishing F-principle for two-layer networks, with extensive controlled
experiments spanning regression, classification, and PDE residual
minimization.  Subsequent work established:
\begin{itemize}
  \item A \emph{general theory}~\cite{Luo2021} for deep networks of
    arbitrary width and depth, analyzing the spectral dynamics of the
    infinite-width neural tangent kernel and yielding closed-form predictions
    for training convergence rates.

  \item An \emph{exactly solvable linear ODE model}~\cite{Luo2022} that
    reproduces F-principle dynamics for two-layer networks, enabling analytic
    predictions of convergence as a function of architecture and learning rate
    without any approximation.

  \item \emph{Phase diagrams}~\cite{PD21} characterizing the distinct
    generalization phases of two-layer ReLU networks at the infinite-width
    limit, providing a rigorous map from initialization and architecture to
    implicit bias and test error.
\end{itemize}

These theoretical results carry concrete engineering consequences for the
design of scientific machine learning methods.  The F-principle predicts
that na\"{i}ve PINN formulations will fail to resolve high-frequency
components of multiscale PDE solutions in stiff regimes---exactly the failure
mode that motivates the APNN architecture.  Conversely, the exact ODE model
of~\cite{Luo2022} provides quantitative guidance on how network depth and
width should scale with the regularity of the target PDE solution, turning
previously heuristic design choices into mathematically principled decisions.

%%--------------------------------------------------------------------------
%%  SECTION 4: INVERSE PROBLEMS
%%--------------------------------------------------------------------------
\section{AI for PDE-Constrained Inverse Problems}

\noindent\textit{Why it matters.}\enspace
Inverse problems governed by PDEs underpin some of the most consequential
computational tasks in modern science and industry.  Full-waveform inversion
(FWI) in geophysics---recovering subsurface wave-speed distributions from
seismic recordings---is the primary computational tool for hydrocarbon
exploration and is increasingly used to monitor underground carbon dioxide
storage, a technology central to achieving climate targets.  Quantitative
photoacoustic and optical tomography in biomedicine requires recovering tissue
optical properties from scattered boundary measurements to achieve diagnostic
precision for cancer detection.  The mathematical challenges are severe: the
forward operators are expensive PDE solves, the problems are severely
ill-posed, and classical approaches based on PDE-constrained optimization or
Markov chain Monte Carlo (MCMC) require thousands to millions of forward
evaluations, making uncertainty quantification computationally prohibitive
at clinical or industrial scale.

\subsection{ODE-Based Diffusion Posterior Sampling}

In~\cite{ODE-DPS2025} we introduced \textbf{ODE-DPS} (ODE-based Diffusion
Posterior Sampling), a framework for solving linear PDE inverse problems
by leveraging pre-trained generative diffusion models as powerful
data-driven physical priors.  The novelty is two-fold.  First, we replace
the conventional stochastic reverse diffusion process with an
adjoint-guided ODE integrator for the posterior score, eliminating the
stochasticity-induced instability of prior diffusion-based samplers and
admitting a rigorous convergence analysis.  Second, we couple the learned
generative prior directly with the analytic likelihood of the physical
forward model, grounding posterior estimation in first principles rather than
heuristics.  Applied to FWI on the OpenFWI benchmark, ODE-DPS achieves
reconstruction quality competitive with fully supervised deep learning models
while requiring an order-of-magnitude fewer labeled training examples---a
critical advantage for real seismic surveys, where ground-truth velocity
models are never available.

\subsection{Bayesian Inversion with Neural Operators}

In~\cite{BINO2025} we proposed \textbf{BINO} (Bayesian Inversion with
Neural Operators), a unified framework that jointly learns the PDE forward
solution operator and performs Bayesian parameter inversion.  A DeepONet
surrogate is trained to map physical parameters to observations; it is then
coupled with a variational Bayes posterior estimator to produce full
uncertainty-quantified posterior distributions over the unknown parameters
at cost orders of magnitude below standard MCMC.  Applied to subdiffusion
models governed by the Caputo fractional derivative---a mathematically
challenging setting where the forward operator is non-local in time---BINO
accurately recovers both the fractional order and diffusion coefficient from
noisy boundary measurements, outperforming classical MCMC and PINN-based
baselines in both posterior accuracy and computational efficiency.

\subsection{Unsupervised Deep Learning for Full-Waveform Inversion}

In~\cite{FWI-DL2025} we developed a fully \emph{unsupervised} deep learning
approach for FWI that requires no paired (velocity model, seismic data)
training examples whatsoever.  The method constructs a physics-constrained
auto-encoder that minimizes the wave-equation residual directly from raw
seismic recordings, producing velocity reconstructions in a zero-shot manner
without any supervision.  This capability is essential for real field
surveys, where labeled data simply do not exist.  Validated on the Marmousi
benchmark---a standard geophysical test case with complex subsurface
structure---the method achieves accuracy competitive with supervised approaches
trained on thousands of labeled instances.

%%--------------------------------------------------------------------------
%%  SECTION 5: FAST SPECTRAL ALGORITHMS
%%--------------------------------------------------------------------------
\section{Fast Spectral Algorithms for Kinetic Equations}

\noindent\textit{Why it matters.}\enspace
Classical fast algorithms play an indispensable dual role in my research
program.  They are the reference solvers that generate ground-truth data
for training and validating all machine learning models, and they are
independent scientific contributions that extend the computational frontier
of kinetic theory.  Without reliable, efficient classical solvers, the
benchmarking of any AI-based method is scientifically unsound.

The direct evaluation of the Boltzmann collision integral at velocity-space
resolution $N^3$ requires $O(N^6)$ operations per time step---completely
intractable for three-dimensional non-equilibrium gas dynamics even on
modern supercomputers.  Fast spectral methods that exploit the convolution
structure of the collision operator in Fourier space reduce this to
$O(N^4 \log N)$ for the elastic Boltzmann operator, but the \emph{inelastic}
Boltzmann operator---governing granular gases, vibrationally excited molecular
gases, and certain plasma-neutral interactions---had not been treated by any
algorithm achieving better than $O(N^6)$ prior to our work.

During my postdoctoral appointment at Purdue University (mentored by
Jingwei~Hu), I developed a \emph{fast spectral method for the inelastic
Boltzmann collision operator}~\cite{HU2019119}, achieving $O(N^2 \log N)$
complexity per velocity mode by uncovering a previously unrecognized
convolution structure in the inelastic collision kernel after a suitable
change of variables in Fourier space.  This is, to our knowledge, the
first sub-cubic algorithm for the inelastic operator.  The method handles
the full non-equilibrium collision dynamics and has been validated against
Direct Simulation Monte Carlo (DSMC) benchmarks for heated granular gas
flows across a broad range of inelasticity parameters.

The implementation is released as open-source software:
\begin{itemize}
  \item \href{https://github.com/mazhengcn/fsm-inelastic-boltzmann}{%
    \texttt{mazhengcn/fsm-inelastic-boltzmann}}: the complete spectral
    solver for the inelastic Boltzmann operator.
  \item \href{https://github.com/mazhengcn/kipack}{\texttt{mazhengcn/kipack}}:
    a high-performance Python package supporting both elastic and inelastic
    Boltzmann operators, with GPU acceleration via JAX/CuPy and a clean
    NumPy interface.
\end{itemize}
These tools are actively used in our group and by external collaborators to
generate training data for APNNs and DeepRTE, closing the loop between fast
classical algorithms and modern machine learning.

Earlier work~\cite{APUQ-Jin2017,APUQ-Jin2018} developed rigorous AP--UQ
methods for uncertain transport and hyperbolic equations via stochastic
Galerkin discretizations, establishing uniform spectral convergence from the
kinetic regime to the diffusion limit in the presence of random inputs.  These
results provided both the mathematical framework and the design intuition that
informed the APNN architecture, and they remain the most complete rigorous
analysis of uncertainty quantification for multiscale kinetic theory.

%%--------------------------------------------------------------------------
%%  SECTION 6: FUTURE DIRECTIONS
%%--------------------------------------------------------------------------
\section{Future Research Directions}

My ongoing and planned research evolves along three interconnected frontiers,
each combining fundamental mathematical questions with high-stakes engineering
applications.

\medskip\noindent\textbf{Foundation models for plasma and fusion confinement.}
Building on the DeepRTE paradigm, I am developing operator-learning
foundation models for the Vlasov--Maxwell system governing fusion plasma
dynamics.  The scientific target is a surrogate model that replaces PIC
simulations for tokamak transport calculations, accurate across the full
kinetic-to-fluid spectrum and evaluable in milliseconds for real-time feedback
control of plasma shape and stability---a prerequisite for commercially viable
fusion energy.  This project is being developed in collaboration with plasma
physicists at SJTU's Institute of Natural Sciences.  A key open challenge is
the preservation of physical invariants (total energy, $L^2$ norm of the
distribution function, entropy inequality) within the operator-learning
architecture, connecting directly to the structure-preserving research thread
below.

\medskip\noindent\textbf{Convergence theory for generative models in inverse problems.}\enspace
Current diffusion-model-based inversion methods, including ODE-DPS, rely on
heuristic justifications for the posterior score approximation step.  I plan
to develop a rigorous error analysis that bounds the Wasserstein distance
between the learned approximate posterior and the true Bayesian posterior as
a function of (i)~the score approximation error, (ii)~the ODE discretization
step size, and (iii)~the forward model accuracy.  This connects F-principle
theory---which governs the spectral convergence of the score network---with
classical Bayesian regularization theory, providing what would be, to our
knowledge, the first end-to-end convergence guarantee for generative inversion
methods.

\medskip\noindent\textbf{Structure-preserving operator learning.}\enspace
Inspired by symplectic integrators and AP schemes in classical numerical
analysis, I aim to build operator-learning architectures that provably
preserve nonlinear conservation laws ($L^p$ mass, energy, entropy inequality)
for Hamiltonian and dissipative PDE systems.  The central technical challenge
is to enforce these constraints as hard architectural invariants rather than
loss-function penalties, since penalty-based approaches are known to degrade
in the stiff multiscale limit---precisely the regime of greatest scientific
interest.  Success here would generalize the APNN design principle to a broad
class of physical constraints applicable to fluid mechanics, plasma physics,
and molecular dynamics simulation.

%%--------------------------------------------------------------------------
%%  SECTION 7: CONCLUSION
%%--------------------------------------------------------------------------
\section{Conclusion}

My research develops a principled, mathematically rigorous approach to AI for
Science that spans from foundational theory to production engineering.  The
unifying insight is that physical structure---the asymptotic limits,
conservation laws, and variational principles that underlie kinetic equations
and their inverses---can and must be encoded as hard constraints in the design
of neural architectures, rather than as optional regularization.  This
produces methods that are not merely empirically successful but provably
correct in the physically relevant limiting regimes.  Equally importantly,
every major methodological contribution in my group is accompanied by an
open-source implementation that substantially lowers the barrier to adoption
for scientists and engineers beyond the computational mathematics community.
I intend to continue building this program at Shanghai Jiao Tong University,
deepening international collaborations, and training computational
mathematicians who can operate fluently at the interface of rigorous analysis
and cutting-edge machine learning.

%%--------------------------------------------------------------------------
%%  BIBLIOGRAPHY
%%--------------------------------------------------------------------------
\vfill
\printbibliography[title={Selected References}]

\end{document}
